% 图7-3：结果独立复核与策略留出验证使用不同的隔离要求。
\begin{tikzpicture}[
  x=1mm,
  y=1mm,
  font=\figurefont\footnotesize,
  text=BookInk,
  record/.style={draw=FigureInput,fill=FigureInput!8,line width=0.8pt,text width=21mm,minimum height=13mm,align=center,inner sep=1.4mm},
  action/.style={draw=FigureModel,fill=FigureModel!8,line width=0.9pt,rounded corners=1.2mm,text width=23mm,minimum height=13mm,align=center,inner sep=1.4mm},
  evidence/.style={draw=BookAmber,fill=BookAmber!8,line width=0.8pt,text width=25mm,minimum height=12mm,align=center,inner sep=1.4mm},
  result/.style={draw=FigureOutput,fill=FigureOutput!6,line width=0.9pt,text width=22mm,minimum height=13mm,align=center,inner sep=1.4mm},
  flow/.style={knowledge arrow,draw=BookTeal,line width=1.0pt,rounded corners=1mm},
  evidence flow/.style={knowledge arrow,draw=BookAmber,line width=0.9pt,rounded corners=1mm},
  boundary/.style={draw=BookMuted,line width=0.8pt,dashed},
  note/.style={font=\figurefont\scriptsize,text=BookMuted,align=center}
]
  \draw[BookRule,line width=0.7pt] (0,37) -- (166,37);
  \node[anchor=north west,font=\figurefont\bfseries] at (1,93) {(a) 已改结果的独立复核：对象可以相同，判断证据需要独立};
  \node[record] (before) at (14,58) {开发对象$x_2$\\修正前记录};
  \node[action] (repair) at (48,58) {依据开发证据\\实施修正};
  \node[record,draw=FigureModel] (after) at (82,58) {开发对象$x_2$\\修正后记录};
  \node[action,draw=BookAmber,fill=BookAmber!8] (review) at (122,58) {未参与修改者\\独立复核};
  \node[result] (verdict) at (153,58) {确认改对／改错\\或无法裁决};
  \node[evidence] (new-evidence) at (122,79) {新证据或盲化材料};

  \draw[flow] (before) -- (repair);
  \draw[flow] (repair) -- (after);
  \draw[flow] (after) -- (review);
  \draw[flow] (review) -- (verdict);
  \draw[evidence flow] (new-evidence) -- (review);

  \node[anchor=north west,font=\figurefont\bfseries] at (1,32) {(b) 总体或策略效果评价：在评价抽样前冻结策略};
  \node[record] (development) at (15,9) {开发审计样本\\已查看问题};
  \node[action] (design) at (50,9) {设计规则、模型\\或接收阈值};
  \node[action,draw=BookInk,fill=BookPaper] (freeze) at (85,9) {冻结修正策略\\和评价口径};
  \node[evidence] (holdout) at (123,9) {独立抽样审计\\比较前后版本};
  \node[result] (effect) at (153,9) {策略效果\\及适用范围};

  \draw[flow] (development) -- (design);
  \draw[flow] (design) -- (freeze);
  \draw[flow] (freeze) -- (holdout);
  \draw[flow] (holdout) -- (effect);
  \draw[boundary] (104,-4) -- (104,22);
  \node[note,anchor=south] at (104,22) {评价边界};
  \node[note,anchor=north,text width=162mm] at (83,-6) {查看评价结果后若继续调整策略，需另取独立证据；开发对象不进入本例的验收总体。};
\end{tikzpicture}
